\documentclass[12pt]{article}
\usepackage[utf8]{inputenc}
\usepackage[T1]{fontenc}
\usepackage{hyperref}
\title{San Miguel French Creole Language Resources for K-12 Curriculum Development}
\author{Compiled by Claude Sonnet 5.5 (claude-sonnet-5-5) for Andr\'as Kornai}
\date{}
\begin{document}
\maketitle

\section{Introduction}
San Miguel French Creole (ISO 639-3 scf), also called Panamanian French Creole or patuá, is a Lesser Antillean French-based creole. Its speakers descend from 19th-century migrants from Saint Lucia and Martinique who came to Panama, mainly to the San Miguel area and Colón, many to work on the Canal. Wikipedia cites 3 speakers (Lee 2017) and ``nearly extinct'' (Snow 2005). Glottolog classifies it as extinct and notes that the Panamanian variety appears not to have been studied. Very little exists, and no school or revitalization materials were found.

\section{Primary Educational Resources}
None found.

\section{Linguistic and Reference Materials}
\begin{itemize}
\item San Miguel Creole (Wikipedia): \url{https://en.wikipedia.org/wiki/San_Miguel_Creole}
\item Glottolog entry for scf: \url{https://glottolog.org/resource/languoid/iso/scf}
\item Verbix Endangered Languages entry: \url{https://docs.verbix.com/EndangeredLanguages/SanMiguelCreoleFrench}
\item Antillean Creole (Wikipedia), the parent variety: \url{https://en.wikipedia.org/wiki/Antillean_Creole}
\end{itemize}
Documentation for the closely related Lesser Antillean creoles (Saint Lucian Kwéyòl, Martinican) is the likeliest source of material for a heritage curriculum, but I did not search for it here and none is listed.

\section{Religious and Community Resources}
None found.

\section{Digital and Audio Resources}
No recordings or corpora of the Panamanian variety were found.

\section{Curriculum Development Context}
There are no known fluent speakers to teach from. A heritage programme would have to be a reclamation effort using the related Antillean creoles as models, together with community memory in the Afro-Antillean communities of Panama. Do not confuse this language with Bocas del Toro Creole, which is English-based, or with the place names San Miguelito and the Bay of San Miguel.

\end{document}